% 张三慧《大学物理学》第4版，第5章《刚体的转动》
% 教材V1 PDF物理页185--189。独立推导候选，未经独立专家审查。
\section*{第5章\quad 刚体的转动：独立解析}

\paragraph{习题 5.1}
圆周运动的切向速率 \(v=R\omega\)，
因此离手角速度
\(\boxed{\omega=25/1.0=25.0\,\mathrm{rad/s}}\)。
若依题意采用\emph{从静止匀角加速}的模型，
转过的角位移
\(\Delta\varphi=1.25(2\pi)=2.5\pi\,\mathrm{rad}\)，
满足 \(\omega^2=2\alpha\Delta\varphi\)，所以
\[
 \boxed{\alpha=\frac{25^2}{5\pi}
 \approx39.79\,\mathrm{rad/s^2}},\qquad
 \boxed{t=\frac{\omega}{\alpha}
 \approx0.628\,\mathrm s}.
\]
原题未显式写 \(\omega_0=0\)；
若允许非零初速，就只有
\(\alpha=(\omega^2-\omega_0^2)/(2\Delta\varphi)\)，
时间不能唯一确定。本题数值依上述起动假设。

\paragraph{习题 5.2}
初、末角速度由 \(n\ \mathrm{r/min}\)
换算为 \(\omega=2\pi n/60\)：
\[
 \boxed{\omega_0=\frac{20\pi}{3}\approx20.94\,\mathrm{rad/s}},\quad
 \boxed{\omega_f=100\pi\approx314.16\,\mathrm{rad/s}}.
\]
恒定角加速度
\[
 \boxed{\alpha=\frac{\omega_f-\omega_0}{7.0}
 =\frac{40\pi}{3}\approx41.89\,\mathrm{rad/s^2}}.
\]
转过角度是平均角速度乘时间：
\[
 \boxed{\Delta\varphi=\frac{\omega_f+\omega_0}{2}t
 =\frac{1120\pi}{3}\,\mathrm{rad}
 \approx1.173\times10^3\,\mathrm{rad}},\quad
 \boxed{N=\Delta\varphi/(2\pi)=560/3\approx186.7\text{ 圈}}.
\]

\paragraph{习题 5.3}
设1900年自转一周周期为 \(T_0\)，
1987年为 \(T_1=T_0+\Delta T\)。
题面给365周累计多费 \(1.14\,\mathrm s\)，
所以
\(\Delta T=1.14/365\approx3.1233\times10^{-3}\,\mathrm s\)。
自转角速度变化
\[
 \Delta\omega=2\pi\!\left(\frac1{T_1}-\frac1{T_0}\right)
 \simeq-\frac{2\pi\Delta T}{T_0^2}.
\]
时间 \(\Delta t=87\times365.25\times86400\,\mathrm s\)，
近似采用 \(T_0\simeq86400\,\mathrm s\)，得
\[
 \boxed{\overline\alpha\simeq
 -\frac{2\pi(1.14/365)}
 {(86400)^2(87\times365.25\times86400)}
 \approx-9.58\times10^{-22}\,\mathrm{rad/s^2}}.
\]
取恒星日代替平太阳日会有很小的系统差异；
此为根据题给平均变慢速率的近似，非实测瞬时角加速度。

\paragraph{习题 5.4}
设地球半径为 \(R_\oplus\)，自转角速度
\(\omega_\oplus\simeq2\pi/(24\,\mathrm h)\)；
纬度 \(\lambda=40^\circ\) 的地点距自转轴
\(r_\perp=R_\oplus\cos\lambda\)。
速率及向心加速度分别为
\[
 \boxed{v=\omega_\oplus R_\oplus\cos40^\circ
 \approx3.55\times10^2\,\mathrm{m/s}},\quad
 \boxed{a=\omega_\oplus^2R_\oplus\cos40^\circ
 \approx2.59\times10^{-2}\,\mathrm{m/s^2}},
\]
其中采用 \(R_\oplus=6.37\times10^6\,\mathrm m\)。
速度沿当地纬线\emph{指向东}；
加速度在垂直自转轴的平面内\emph{指向自转轴}，
不能直接说“指向地心”（两者在非赤道处不同）。

\paragraph{习题 5.5}
把氧原子视作坐标原点，
两氢原子位于
\((d\cos(\theta/2),\pm d\sin(\theta/2))\)。
关于水平 \(AA'\) 轴、竖直 \(BB'\) 轴的转动惯量分别为
\[
 J_{AA'}=2m_Hd^2\sin^2(\theta/2),\qquad
 J_{BB'}=2m_Hd^2\cos^2(\theta/2).
\]
故
\[
 \boxed{d=
 \sqrt{\frac{J_{AA'}+J_{BB'}}{2m_H}}
 \approx9.58\times10^{-11}\,\mathrm m},
\qquad
 \boxed{\theta=2\arctan\sqrt{\frac{J_{AA'}}{J_{BB'}}}
 \approx104.91^\circ}.
\]
原图 \(\theta\) 是两条OH键的完整夹角，
单条OH键与水平 \(AA'\) 正轴的夹角为
\(\theta/2\approx52.46^\circ\)。
辅导书PDF物理页81也以 \(\sin^2(\theta/2)\)
及 \(\cos^2(\theta/2)\) 计算，二者一致。
氧原子在双轴交点，其平动位置不贡献转动惯量。

\paragraph{习题 5.6}
均匀球壳质量等于60个碳原子之和，
\(M\simeq60(12u)=720u
 \approx1.196\times10^{-24}\,\mathrm{kg}\)。
按半径 \(R=D/2\) 的均匀薄球面，
绕任一直径的转动惯量
\[
 \boxed{J=\frac23MR^2=\frac16MD^2}.
\]
\textbf{必须保留扫描源疑点：}原题PDF物理页186
印为 \(D=71\,\mathrm{nm}\)，
由此严格按字面得到
\[
 \boxed{J_{\rm printed}\approx1.00\times10^{-39}
 \,\mathrm{kg\,m^2}},\qquad
 f_{\rm printed}=\frac1{2\pi}
 \sqrt{\frac{2E_k}{J_{\rm printed}}}
 \approx\boxed{5.60\times10^8\,\mathrm{Hz}}.
\]
但这明显大于常见富勒烯直径的两个数量级。
辅导书PDF物理页81同样将直径印为
\(71\,\mathrm{nm}\)，且采用
\(E_k=6.12\times10^{-21}\,\mathrm J\)，
与教材本卷第186页给出的
\(6.21\times10^{-21}\,\mathrm J\) 有数值分歧。
辅导书据此得到约 \(5.54\times10^8\,\mathrm{Hz}\)；
按照教材的 \(6.21\times10^{-21}\,\mathrm J\)
计算得到上面的 \(5.60\times10^8\,\mathrm{Hz}\)。
如果出于物理尺寸考虑猜测原意是
\(D=0.71\,\mathrm{nm}\)，则
\[
 \boxed{J_{\rm physical}\approx1.00\times10^{-43}
 \,\mathrm{kg\,m^2}},\quad
 \boxed{f_{\rm physical}\approx5.60\times10^{10}\,\mathrm{Hz}}.
\]
两组结果分别对应两个不同输入条件，
不能直接改掉原题，也不能在题目未核验时
只留下看起来合理的一组数值。

\paragraph{习题 5.7}
氧分子中两氧原子对质心的距离均为
\(d/2\)，其中 \(d=1.21\times10^{-10}\,\mathrm m\)。
把原子看成质点：
\[
 \boxed{J=2m_O(d/2)^2=\frac12m_Od^2
 \approx1.947\times10^{-46}\,\mathrm{kg\,m^2}}.
\]
由 \(E_k=\tfrac12J\omega^2\) 求
\(\omega=\sqrt{2E_k/J}\)，所以
\[
 \boxed{T=\frac{2\pi}{\omega}
 =2\pi\sqrt{\frac{J}{2E_k}}
 \approx1.37\times10^{-12}\,\mathrm s}.
\]
为经典刚性双原子模型，不包括量子转动能级。

\paragraph{习题 5.8}
两个铁球质量各为 \(m\)，半径各为 \(R\)；
杆在球面间长 \(\ell\)，故球心离杆中心
\(a=R+\ell/2\)。
每个均匀实心球对过球心的直径
转动惯量为 \(\tfrac25mR^2\)。
平行轴定理给出垂直连杆并过中点的总转动惯量
\[
 \boxed{J_\perp
 =2\left(\tfrac25mR^2+ma^2\right)
 =\frac45mR^2+2m(R+\ell/2)^2}.
\]
对沿两球心连线的轴无平移轴项，
\[
 \boxed{J_\parallel=2(\tfrac25mR^2)=\frac45mR^2}.
\]
若把杆长 \(\ell\) 错当作球心距离，结果会偏小；
这里依据原图5.33的杆段位置。

\paragraph{习题 5.9}
将两指针近似为以一端绕表盘中心转动的均匀细杆，
\(J=\tfrac13mL^2\)。
分针：
\[
 J_m=\tfrac13(100)(4.50)^2=675\,\mathrm{kg\,m^2},
 \quad\omega_m=\frac{2\pi}{3600}\,\mathrm{rad/s}.
\]
时针：
\[
 J_h=\tfrac13(60.0)(2.70)^2=145.8\,\mathrm{kg\,m^2},
 \quad\omega_h=\frac{2\pi}{12(3600)}\,\mathrm{rad/s}.
\]
因此各自角动量和转动动能为
\[
 \boxed{L_m=J_m\omega_m\approx1.18\,\mathrm{kg\,m^2/s}},
 \quad \boxed{K_m=\tfrac12J_m\omega_m^2
 \approx1.03\times10^{-3}\,\mathrm J},
\]
\[
 \boxed{L_h=J_h\omega_h\approx2.12\times10^{-2}
 \,\mathrm{kg\,m^2/s}},\quad
 \boxed{K_h=\tfrac12J_h\omega_h^2
 \approx1.54\times10^{-6}\,\mathrm J}.
\]
对表盘观察，其转向一致，矢量角动量同向。

\paragraph{习题 5.10}
设原薄圆盘质量为 \(M_0\)，
挖去半径 \(R/2\) 的小盘质量
\(M_h=M_0/4\)，余板质量
\(m=3M_0/4\)。
故 \(M_0=4m/3,\ M_h=m/3\)。
由平行轴定理，原中心处挖去部分的转动惯量
\[
 J_h=\tfrac12M_h(R/2)^2+M_h(R/2)^2
 =\tfrac18mR^2.
\]
原完整圆盘关于原中心
\(J_0=\tfrac12M_0R^2=\tfrac23mR^2\)。
余板惯量为
\[
 \boxed{J=J_0-J_h=\frac{13}{24}mR^2}.
\]
检查：小盘面积与原盘之比为 \(1/4\)，
不是半径之比 \(1/2\)。

\paragraph{习题 5.11}
取悬挂物 \(m_1\) 向下为正，
桌上物 \(m_2\) 向滑轮移动为正，
若有滑动摩擦 \(f_k=\mu_km_2g\)，
两物共用加速度 \(a\)。
设悬挂侧张力 \(T_1\)、桌面侧张力 \(T_2\)，
有
\[
 m_1g-T_1=m_1a,\qquad
 T_2-\mu_km_2g=m_2a.
\]
滑轮为均匀圆盘 \(J=\tfrac12mr^2\)，
且绳与轮不打滑：
\((T_1-T_2)r=J(a/r)\)。
联立得
\[
 \boxed{a=\frac{(m_1-\mu_km_2)g}
 {m_1+m_2+m/2}},\quad
 \boxed{T_1=m_1(g-a)},\quad
 \boxed{T_2=m_2a+\mu_km_2g}.
\]
该式在 \(m_1>\mu_km_2\) 时
对应题设的向下运动，且假定已进入滑动状态；
静摩擦启动阈值不同，不能据此证明必会启动。

\paragraph{习题 5.12}
米尺长 \(L=1.00\,\mathrm m\)，
均匀质心在 \(50\,\mathrm{cm}\) 处，
墙上转轴在 \(60\,\mathrm{cm}\) 刻度，
质心偏离轴 \(b=0.10\,\mathrm m\)。
由平行轴定理
\[
 J_p=m(L^2/12+b^2).
\]
水平初态受重力矩大小 \(mgb\)，
故
\[
 \boxed{\alpha_0=\frac{mgb}{J_p}
 =\frac{gb}{L^2/12+b^2}
 =10.5\,\mathrm{rad/s^2}}.
\]
竖直到最低点时质心下落 \(b\)，
机械能守恒
\(\tfrac12J_p\omega^2=mgb\)，得
\[
 \boxed{\omega=\sqrt{\frac{2gb}{L^2/12+b^2}}
 \approx4.58\,\mathrm{rad/s}}.
\]
忽略轴摩擦与空气阻力。

\paragraph{习题 5.13}
定轴刚体角速度对全部质量微元相同为 \(\omega\)。
距轴垂直距离为 \(r_i\) 的微元
线速率 \(v_i=r_i\omega\)，
总动能
\[
 K=\sum_i\frac12\Delta m_i v_i^2
 =\frac12\omega^2\sum_i\Delta m_i r_i^2.
\]
根据转动惯量定义
\(J=\sum_i\Delta m_i r_i^2\)，
连续极限 \(J=\int r_\perp^2\,dm\)，
故
\[
 \boxed{K=\frac12J\omega^2}.
\]
物体绕定轴转动时上述表达式成立；
有质心平移的刚体一般还需加平动动能。

\paragraph{习题 5.14}
在唱片与转盘存在相对滑动时，
法向压力可按均匀盘面分布。
在半径 \(r\) 处的环带质量
\(dm=(2mr/R^2)\,dr\)，
该环带动摩擦力微元
\(dF=\mu_kg\,dm\)，
其力矩 \(d\tau=r\,dF\)。
积分得总摩擦矩
\[
 \tau=\int_0^Rr\mu_kg\frac{2mr}{R^2}\,dr
 =\frac23\mu_kmgR.
\]
唱片为均匀圆盘 \(J=\tfrac12mR^2\)，
角加速度
\[
 \alpha=\tau/J=\frac{4\mu_kg}{3R}.
\]
由静止加速到转盘角速度 \(\omega\) 需
\[
 \boxed{t=\frac{3R\omega}{4\mu_kg}}.
\]
理想转盘维持 \(\omega\) 不变，
电动机做功
\[
 W_{\rm motor}=\tau\,\omega\,t
 =J\omega^2=\boxed{\frac12mR^2\omega^2}.
\]
唱片增加的转动动能为
\[
 \boxed{\Delta K=\frac12J\omega^2
 =\frac14mR^2\omega^2}.
\]
其余 \(\Delta K\) 量级的电机功转变为
两表面相对滑动的摩擦热；
不能把电机功与唱片动能误认为全相等。

\paragraph{习题 5.15}
转椅轴上外力矩可忽略，
因此收臂过程角动量守恒
\(J_1\omega_1=J_2\omega_2\)。
故
\[
 \boxed{n_2=n_1\frac{J_1}{J_2}
 =15(2/0.80)=37.5\,\mathrm{r/min}}.
\]
机械能并不守恒，收臂的人必须克服
由旋转惯性产生的拉力做正功：
\[
 W_{\rm person}=\Delta K
 =\frac12J_2\omega_2^2-\frac12J_1\omega_1^2
 =\frac{J_1\omega_1^2}{2}
 \left(\frac{J_1}{J_2}-1\right)
 \approx\boxed{3.70\,\mathrm J}.
\]
能量来自人体内部化学能；
角动量守恒与机械能变化并不矛盾。


\paragraph{习题 5.16}
子弹质量 \(m=0.0080\,\mathrm{kg}\)，
速度 \(v_0=200\,\mathrm{m/s}\)，
入射点距转轴 \(d=3L/4=0.30\,\mathrm m\)。
碰撞时间极短，转轴冲量对该轴的力矩为零，
故碰撞前后关于固定轴的角动量守恒：
\[
 mv_0d=J_{\rm tot}\omega_0,\quad
 J_{\rm tot}=\frac13ML^2+md^2.
\]
代入 \(L=0.40\,\mathrm m,M=1.0\,\mathrm{kg}\)，
\[
 J_{\rm tot}\approx0.0540533\,\mathrm{kg\,m^2},
 \qquad
 \boxed{\omega_0=\frac{mv_0d}{J_{\rm tot}}
 \approx8.88\,\mathrm{rad/s}}.
\]
子弹嵌入造成非弹性损失，
碰后\emph{立即开始摆动阶段}机械能才守恒。
摆到最大偏角 \(\theta\) 时角速度为零，
杆质心提高 \(L(1-\cos\theta)/2\)，
子弹提高 \(d(1-\cos\theta)\)，因此
\[
 \frac12J_{\rm tot}\omega_0^2
 =g\left(\frac{ML}{2}+md\right)(1-\cos\theta).
\]
解得
\[
 \boxed{\theta_{\max}=
 \arccos\left[1-
 \frac{(mv_0d)^2}{
 2J_{\rm tot}g(ML/2+md)}\right]
 \approx94.3^\circ}.
\]
若右侧 \(1-\cos\theta>2\) 则可能完成整圈，
但本组参数约为 \(1.074<2\)，仍有转折点。

\paragraph{习题 5.17}
转台转动惯量 \(J=1200\,\mathrm{kg\,m^2}\)，
周期 \(10\,\mathrm s\)，
初角速度 \(\omega_0=2\pi/10\)。
人从中心移动到半径 \(r=2\,\mathrm m\)，
人的质量为 \(M=80\,\mathrm{kg}\)，
若忽略竖直轴的外力矩，
\[
 J\omega_0=(J+Mr^2)\omega_f.
\]
因此
\[
 \boxed{\omega_f=
 \frac{1200}{1200+80(2)^2}\frac{2\pi}{10}
 \approx0.496\,\mathrm{rad/s}}.
\]
转速下降对应转动惯量增加；
人在转台上走动可由内部作用改变系统动能，
但不改变总轴向角动量。

\paragraph{习题 5.18}
两车等质量 \(m=1200\,\mathrm{kg}\)，
速度大小 \(v=72/3.6=20.0\,\mathrm{m/s}\)，
迎面、平行行驶方向的质心轨迹间距
\(b=0.50\,\mathrm m\)。
相对于两车总体质心，
碰前总动量为零，但总角动量不为零：
\[
 L_{\rm before}=mvb
 =1200(20)(0.50)=1.20\times10^4\,\mathrm{kg\,m^2/s}.
\]
由短时碰撞外力矩可忽略，
且碰后残体 \(J=2500\,\mathrm{kg\,m^2}\)，得
\[
 \boxed{|\omega|=\frac{|L|}{J}=4.80\,\mathrm{rad/s}}.
\]
初动能为
\(K_i=2(\tfrac12mv^2)=4.80\times10^5\,\mathrm J\)，
碰后仅旋转动能
\(K_f=\tfrac12J\omega^2
 =2.88\times10^4\,\mathrm J\)。
非弹性碰撞损失机械能
\[
 \boxed{E_{\rm lost}=K_i-K_f
 =4.512\times10^5\,\mathrm J}.
\]
力学动量守恒不意味着动能守恒。

\paragraph{习题 5.19}
(1) 宇航员在舱内壁沿圆周跑动，
要在自身共转参考系中有大小 \(g\) 的向心加速度，
\(v^2/r=g\)，\(r=2.5\,\mathrm m\)，故
\[
 \boxed{v=\sqrt{gr}=\sqrt{9.8\times2.5}
 \approx4.95\,\mathrm{m/s}}.
\]
(2) 三人总转动惯量
\(J_a=3m_ar^2=3(70)(2.5)^2
 =1312.5\,\mathrm{kg\,m^2}\)，
船体 \(J_s=3\times10^5\,\mathrm{kg\,m^2}\)。
若船最初静止，三人沿船壁相对船的
角速度为 \(\omega_{\rm rel}=v/r\)，
令船体反向转速为 \(\Omega\)，
人对惯性系角速度为
\(\Omega+\omega_{\rm rel}\)。
总角动量守恒给出
\[
 0=J_s\Omega+J_a(\Omega+\omega_{\rm rel}),
\quad
 \boxed{\Omega=-\frac{J_a}{J_s+J_a}\frac vr
 \approx-8.62\times10^{-3}\,\mathrm{rad/s}}.
\]
(3) 船转过 \(30^\circ\) 需时间
\(t=(\pi/6)/|\Omega|\approx60.7\,\mathrm s\)。
三人在船壁上各跑过圈数
\[
 \boxed{N=\frac{vt}{2\pi r}
 =\frac{J_s+J_a}{J_a}\frac{1}{12}
 \approx19.1\text{ 圈}}.
\]
本解把 \(v\) 理解为相对舱壁的跑速。
若把第一问的 \(v\) 解释为对惯性系速度，
分母应改为 \(J_s\)，结果略有差别；
必须明确参考系，不将两者混为一谈。

\paragraph{习题 5.20}
把太阳视作均匀实心球，
则关于自转轴的
\(J_\odot=(2/5)M_\odot R_\odot^2\)。
假设太阳自转周期 \(T_\odot=26\,\mathrm d\)，
以本书数值表量级
\(M_\odot\approx1.989\times10^{30}\,\mathrm{kg}\)、
\(R_\odot\approx6.96\times10^8\,\mathrm m\)，
得到
\[
 \boxed{L_\odot=J_\odot\omega_\odot
 =\frac25M_\odot R_\odot^2
 \frac{2\pi}{T_\odot}
 \approx1.08\times10^{42}\,\mathrm{kg\,m^2/s}}.
\]
题给太阳系总角动量
\(L_{\rm SS}=3.3\times10^{43}\,\mathrm{kg\,m^2/s}\)，
故
\[
 \boxed{100\%\times L_\odot/L_{\rm SS}\approx3.26\%}.
\]
若太阳差速自转、密度不均匀，
结果与均匀实心球模型不同。

\paragraph{习题 5.21}
周期 \(T=0.033\,\mathrm s\)，
周期增长速率
\(dT/dt=1.26\times10^{-5}\,\mathrm{s/a}\)。
令一年为 \(365.25\times86400\,\mathrm s\)，
由 \(\omega=2\pi/T\) 得
\[
 \boxed{\alpha=\frac{d\omega}{dt}
 =-\frac{2\pi}{T^2}\frac{dT}{dt}
 \approx-2.30\times10^{-9}\,\mathrm{rad/s^2}}.
\]
把星体近似为质量
\(M=1.5\times10^{30}\,\mathrm{kg}\)、
半径 \(R=10^4\,\mathrm m\) 的均匀球，
\(J=(2/5)MR^2
 =6.0\times10^{37}\,\mathrm{kg\,m^2}\)。
自转动能减小率
\[
 \boxed{-\frac{dK}{dt}
 =-J\omega\alpha
 =J\frac{2\pi}{T}|\alpha|
 \approx2.63\times10^{31}\,\mathrm{J/s}}.
\]
题(3)明确假定\emph{这一能量变化率不变}，
并非假定角减速度不变。
初始转动动能
\(K_0=\tfrac12J\omega^2\)，
能量耗尽时间
\[
 \boxed{t_{\rm stop}=\frac{K_0}{|dK/dt|}
 =\frac{\omega}{2|\alpha|}
 =\frac{T}{2(dT/dt)}
 \approx1.31\times10^3\,\mathrm a}.
\]
如果错误采用恒角减速度，
会得到上式两倍的时间。
真实脉冲星减速律不会在全部寿命内严格保持恒定功率。

\paragraph{习题 5.22}
将地球视作转动惯量
\(J=0.33M_\oplus R_\oplus^2\) 的刚体，
周期约 \(T=86400\,\mathrm s\)。
取 \(M_\oplus=5.972\times10^{24}\,\mathrm{kg}\)、
\(R_\oplus=6.371\times10^6\,\mathrm m\)。
自转动能
\[
 \boxed{K=\frac12J(2\pi/T)^2
 \approx2.12\times10^{29}\,\mathrm J}.
\]
给定每 \(10^6\) 年周期增大 \(16\,\mathrm s\)，
因此
\[
 \dot T=\frac{16}
 {10^6(365.25)(86400)}
 \quad(\mathrm{s/s}).
\]
由 \(\dot\omega=-2\pi\dot T/T^2\) 得
\[
 \boxed{P_{\rm loss}=-\dot K
 =J\omega|\dot\omega|
 \approx2.48\times10^{12}\,\mathrm W},
\]
\[
 \boxed{E_{\rm year}=P_{\rm loss}(365.25)(86400)
 \approx7.83\times10^{19}\,\mathrm J
 \approx2.18\times10^{13}\,\mathrm{kW\,h}}.
\]
与我国2004年发电量的比较，
配套辅导书PDF物理页88采用能量
\(E_{2004}=7.3\times10^{18}\,\mathrm J\)
（该页题面括号把单位印为「kW」，
但其计算使用的是能量单位 \(\mathrm J\)，
为源书量纲疑点）。
据此
\[
 \boxed{\frac{E_{\rm year}}{E_{2004}}
 =\frac{7.83\times10^{19}}{7.3\times10^{18}}
 \approx10.7\text{ 倍}\approx11\text{ 倍}}.
\]
采用其他实际2004年统计口径应相应调整比值，
不能把印刷的 \(\mathrm{kW}\) 当作发电能量。
作用在地球上的平均摩擦力矩大小
\[
 \boxed{|\tau|=J|\dot\omega|
 \approx3.41\times10^{16}\,\mathrm{N\,m}}.
\]
这里是题设长期平均周期漂移的数量级模型。

\paragraph{习题 5.23}
忽略质量散失，塌缩前后太阳对自转轴
无外部净力矩近似，角动量守恒：
\(J_i\omega_i=J_f\omega_f\)。
均匀球的 \(J\propto MR^2\)，
所以
\[
 \boxed{T_f=T_i\left(\frac{R_f}{R_i}\right)^2}.
\]
取 \(T_i=26\,\mathrm d\)、
白矮星最终半径按地球半径
\(R_f=6.37\times10^6\,\mathrm m\)，
当前太阳半径
\(R_i=6.96\times10^8\,\mathrm m\)，得到
\[
 \boxed{T_f\approx188\,\mathrm s
 \approx3.14\,\mathrm{min}}.
\]
实际白矮星形成可能丢失质量和角动量，
本结果仅是题设无质量散失的理想估算。

\paragraph{习题 5.24}
设圆盘质心为 \(C\)，与固定地面
瞬时接触点为 \(P\)，盘上任意物质点为 \(A\)。
刚体速度叠加：
\[
 \boldsymbol v_A=\boldsymbol v_C
 +\boldsymbol\omega\times\overrightarrow{CA}.
\]
无滑动滚动意味着
\(\boldsymbol v_P=\boldsymbol0\)，
于是
\[
 \boldsymbol v_C=-\boldsymbol\omega
 \times\overrightarrow{CP}
 =\boldsymbol\omega\times\overrightarrow{PC}.
\]
代回：
\[
 \boldsymbol v_A
 =\boldsymbol\omega\times
 [\overrightarrow{PC}+\overrightarrow{CA}]
 =\boxed{\boldsymbol\omega\times\overrightarrow{PA}}.
\]
因此以 \(P\) 为瞬时速度零点的瞬时转动中心
是几何接触点；这不表示 \(P\)
在整个运动中一直是同一个空间固定的轴，
也不意味着接触点的加速度必为零。

\paragraph{习题 5.25}
轮轴外半径 \(R_1=1.00\,\mathrm m\)，
内缆绕半径 \(R_2=0.60\,\mathrm m\)，
质量 \(M=1000\,\mathrm{kg}\)，
质心转动惯量 \(J_C=300\,\mathrm{kg\,m^2}\)。
电缆在内轴下方受向右拉力 \(F=2000\,\mathrm N\)。
以与地面接触点 \(P\) 为瞬时转动中心，
\(J_P=J_C+MR_1^2=1300\,\mathrm{kg\,m^2}\)。
拉力作用线高于接地点
\(R_1-R_2=0.40\,\mathrm m\)，
使轮轴顺时针加速滚动，
故
\[
 \boxed{|\alpha|=\frac{F(R_1-R_2)}{J_C+MR_1^2}
 =\frac{2000(0.40)}{1300}
 \approx0.615\,\mathrm{rad/s^2}}.
\]
质心向右加速度大小
\[
 \boxed{a_C=R_1|\alpha|
 \approx0.615\,\mathrm{m/s^2}}.
\]
设地面摩擦水平分量
\(f\) 向右为正，质心动量定理
\(F+f=Ma_C\)：
\[
 \boxed{f=Ma_C-F\approx-1.385\times10^3\,\mathrm N}.
\]
即静摩擦实际向左，大小 \(1.38\,\mathrm{kN}\)。
滚动无滑动可行的前提是地面最大静摩擦力
能提供这一数值；题设未给摩擦系数，
故这是运动满足无滑动时所\emph{需}的静摩擦。

\paragraph{习题 5.26}
把「悠悠球」简化为均匀实心圆柱，
质量 \(m\)、半径 \(R\)，
放出的线绕在圆柱外缘，不打滑。
关于水平轴
\(J=\frac12mR^2\)。
(1) 质心向下加速度 \(a\)，线张力 \(T\) 向上，
\[
 mg-T=ma,\quad TR=J\alpha,\quad a=R\alpha.
\]
联立得
\[
 \boxed{a=\frac{g}{1+J/(mR^2)}=\frac23g},
 \qquad
 \boxed{T=mg-ma=\frac13mg}.
\]
(2) 若圆柱体质心要保持恒定高度，
则 \(a_C=0\)，所以
\(\boxed{T=mg}\)，
但是仍可有非零角加速度
\[
 \boxed{\alpha=\frac{TR}{J}=\frac{2g}{R}}.
\]
(3) 若线端向上拉力 \(F=T=2mg\)，
则圆柱体质心竖直加速度为
\[
 \boxed{a_C=\frac{T-mg}{m}=g\quad\text{向上}},
\]
绕质心角加速度大小
\(\boxed{\alpha=TR/J=4g/R}\)。
对于沿圆柱切线释放的线端，
无滑动运动学意味着线端加速度
为质心向上加速度与外缘切向加速度之和
（这两者在题设拉线方向上同向），
故
\[
 \boxed{a_{\rm hand}=a_C+R\alpha
 =g+4g=5g\quad\text{向上}}.
\]
如线缠绕的\emph{轴半径}不同于外圆柱半径，
必须引入独立绕线半径，不能沿用此数值。

\paragraph{习题 5.27}
地球自旋轴与轨道面法线的恒定夹角为
\(\theta=23.5^\circ\)。
自转角动量矢量 \(\boldsymbol L\)
以进动角速度
\(\Omega=2\pi/(26000\,\mathrm a)\)
绕轨道法线旋转，
大小 \(L=J\omega\) 近似不变。
矢量在圆锥面上运动，故
\[
 \left|\frac{d\boldsymbol L}{dt}\right|
 =\Omega L\sin\theta
 =\frac{2\pi}{26000(365.25)(86400)}
  \left(\frac{2\pi J}{86400}\right)
 \sin23.5^\circ.
\]
代入 \(J=8.05\times10^{37}\,\mathrm{kg\,m^2}\)，
可得
\[
 \boxed{|d\boldsymbol L/dt|
 \approx1.79\times10^{22}\,\mathrm{N\,m}}.
\]
角动量定理
\(\boldsymbol\tau_{\rm external}
 =d\boldsymbol L/dt\)；
所以太阳和月亮对地球的\emph{合力矩大小}
亦约为
\(\boxed{1.79\times10^{22}\,\mathrm{N\,m}}\)。
不能错误地漏掉圆锥进动的 \(\sin\theta\) 因子。

\paragraph{习题 5.28}
均匀实心圆盘质量
\(M=50.0\,\mathrm t=5.00\times10^4\,\mathrm{kg}\)，
半径 \(R=2.00\,\mathrm m\)，
惯量
\[
 J=\frac12MR^2=1.00\times10^5\,\mathrm{kg\,m^2}.
\]
额定角速度
\(\omega=800(2\pi/60)\approx83.78\,\mathrm{rad/s}\)，
从静止转至额定速度需动能
\(K=\tfrac12J\omega^2\)。
恒功率 \(P=7.46\times10^4\,\mathrm W\) 时
\[
 \boxed{t=\frac{J\omega^2}{2P}
 \approx4.70\times10^3\,\mathrm s
 \approx78.4\,\mathrm{min}}.
\]
盘的自旋角动量大小
\(L=J\omega\approx8.38\times10^6\,\mathrm{kg\,m^2/s}\)。
船体使自旋轴以
\(\Omega=1^\circ/\mathrm s=\pi/180\,\mathrm{rad/s}\)
垂直于 \(\boldsymbol L\) 进动时，
所需进动力矩大小
\[
 \boxed{|\tau|=|\boldsymbol\Omega\times\boldsymbol L|
 =\Omega L\approx1.46\times10^5\,\mathrm{N\,m}}.
\]
船体受回转体的反作用力矩大小相同、方向相反。
左右倾方向还须指定旋盘自旋方向
与船体进动的矢量方向；
没有这两个箭头时不能唯一给出
「左倾/右倾」符号。
